\begin{definition}[Iterator Relation]
\label{def:iterator-relation}
Let $(P,S,1)$ be a Peano system, and let $(W,c,g)$ be iterator data for it. A
relation $R\subseteq P\times W$ is called an \emph{iterator relation} for
$(W,c,g)$ exactly when
\[
(1,c)\in R
\]
and
\[
\forall n\in P\;\forall w\in W\;
\bigl((n,w)\in R\Longrightarrow (S(n),g(w))\in R\bigr).
\]
\end{definition}
\end{definitionbox}

\begin{remark*}[Standard quantified statement]
\[
\begin{aligned}
R \text{ is an iterator relation for } (W,c,g)
\Longleftrightarrow\;&
R\subseteq P\times W
\land (1,c)\in R\\
&\land
\forall n\in P\;\forall w\in W\;
\bigl((n,w)\in R\Longrightarrow (S(n),g(w))\in R\bigr).
\end{aligned}
\]
\end{remark*}

\begin{remark*}[Predicate reading]
\[
\begin{aligned}
\operatorname{IteratorRelation}(R,(P,S,1),W,c,g)
\Longleftrightarrow\;&
\operatorname{Relation}(R,P,W)
\land (1,c)\in R\\
&\land
\forall n\in P\;\forall w\in W\;
\bigl((n,w)\in R\Longrightarrow (S(n),g(w))\in R\bigr).
\end{aligned}
\]
\end{remark*}

\begin{remark*}[Negated quantified statement]
\[
\begin{aligned}
R \text{ is not an iterator relation for } (W,c,g)
\Longleftrightarrow\;&
R\not\subseteq P\times W
\lor (1,c)\notin R\\
&\lor \exists n\in P\;\exists w\in W\;
\bigl((n,w)\in R\land (S(n),g(w))\notin R\bigr).
\end{aligned}
\]
\end{remark*}

\begin{remark*}[Failure modes]
\begin{description}
  \item[Exposition.]
  An iterator relation fails if it is not a relation inside $P\times W$, if
  it omits the required base pair, or if it is not closed under the iterator
  step.
  \item[Invalid relation.]
  \[
  \neg\operatorname{IteratorRelation}(R,(P,S,1),W,c,g).
  \]
  \[
  \begin{aligned}
  R \text{ is not an iterator relation for } (W,c,g)
  \Longleftrightarrow\;&
  \underbrace{R\not\subseteq P\times W}_{\text{wrong ambient relation}}\\
  &\lor \underbrace{(1,c)\notin R}_{\text{base pair missing}}\\
  &\lor \underbrace{\exists n\in P\;\exists w\in W\;
  \bigl((n,w)\in R\land (S(n),g(w))\notin R\bigr)}_{\text{step closure fails}}.
  \end{aligned}
  \]
\end{description}
\end{remark*}

\begin{remark*}[Interpretation]
An iterator relation is any graph-like set of possible stage-value pairs that
contains the required first pair and is closed under the recursive step.
\end{remark*}

\begin{dependencies}
\begin{itemize}
  \item \hyperref[def:iterator-data-on-peano-system]{Iterator Data}
\end{itemize}
\end{dependencies}

